\chapter{Reading Public Earnings Releases, Tech 10-Qs, and AI Startup S-1s}
\label{ch:filings-and-profitability}
\chapterline

In Chapter~\ref{ch:valuation-multiples}, quantities such as Free Cash Flow, Enterprise Value, EBITDA, and Return on Invested Capital entered our equations as clean algebraic variables. In practice, nobody hands you clean variables. Public companies publish twenty-page quarterly earnings releases packed with ``Non-GAAP'' adjustments alongside hundred-page regulatory filings; pre-IPO technology startups file S-1 prospectuses blending positive unit cohorts with multibillion-dollar accounting losses; private business brokers email QuickBooks exports padded with questionable add-backs.

To evaluate any operating company in thirty minutes without getting misled by headline adjustments, an engineer needs a universal diagnostic toolkit first -- segment disaggregation, variable cost of revenue versus fixed operating overhead, equity compensation dilution, capital-expenditure-to-depreciation timing lags, unit contribution versus recurring R\&D step-costs, cohort retention math, and owner-compensation normalization. In each section below, we establish those general accounting and economic laws first and then apply them to an explicit illustrative case study: a public hyperscaler quarterly release (Alphabet reporting archetype), a frontier AI lab or cloud S-1 (Anthropic-style serving versus pre-training archetype), and a mid-sized private operating business.

\begin{historybox}
Before 1933, corporations selling stock on U.S. exchanges operated under a patchwork of state ``blue sky'' laws with no uniform federal accounting rules. Industrial trusts routinely published balance sheets without income statements, concealed debt inside unconsolidated holding-company pyramids, or booked unrealized asset write-ups as operating profit. After the 1929 market crash wiped out over $80\%$ of U.S. equity market capitalization and the Senate's Pecora Commission exposed widespread insider self-dealing, Congress enacted two statutes built around Louis Brandeis's principle of mandatory public disclosure rather than government merit approval.

The \textbf{Securities Act of 1933} governs the \emph{primary} offering of securities to the public: before selling shares in an Initial Public Offering (IPO), an issuer must file a registration statement -- \textbf{Form S-1} (or \textbf{Form F-1} for foreign private issuers) -- containing audited financial statements, unit economics, cap-table dilution, executive pay, and explicit Risk Factors, with strict statutory liability for material omissions. The \textbf{Securities Exchange Act of 1934} established the \textbf{Securities and Exchange Commission (SEC)} and mandated ongoing secondary-market reporting: \textbf{Form 10-K} (annual report with independently audited statements and full footnotes), \textbf{Form 10-Q} (quarterly unaudited statements reviewed by auditors), and \textbf{Form 8-K} (current report filed within four business days of material events such as acquisitions, debt financings, auditor resignations, or quarterly earnings releases). While statutory filings must follow \textbf{GAAP} (U.S. Generally Accepted Accounting Principles) or \textbf{IFRS} (International Financial Reporting Standards), modern press releases often headline management's custom ``Non-GAAP'' or ``Adjusted'' figures -- which SEC Regulation~G requires to be reconciled line by line back to GAAP.
\end{historybox}

\section{Dissecting Public Quarterly Releases and 10-Q/10-K Filings}

When a public company reports quarterly results, skip the press-release narrative bullets and open four tables immediately: the Consolidated Statement of Income, the Disaggregated Revenue and Segment Footnote, the Consolidated Balance Sheet, and the Consolidated Statement of Cash Flows. Four universal accounting principles govern how you translate those tables into economic cash power.

\subsection{Core principles: segments, distribution COGS, SBC dilution, and CapEx waves}

\begin{enumerate}
  \item \textbf{Segment disaggregation, operating leverage, and sum-of-the-parts (ASC~280 / IFRS~8).}
  Consolidated totals average together divisions with vastly different marginal costs and capital intensities. Accounting standards require public companies to report revenue and operating profit along the same segment lines reviewed by the Chief Operating Decision Maker. Two diagnostic checks belong in every segment review:
  \begin{itemize}
    \item \emph{Incremental operating margin:} For a business scaling over a large fixed-cost infrastructure base, compute $\frac{\Delta \text{Operating Income}}{\Delta \text{Revenue}}$ year over year. When fixed costs are already paid for, incremental operating margins can exceed $50\%$ even when average segment margin is only $15\%$.
    \item \emph{Isolating venture-stage losses (sum-of-the-parts):} If Segment~A earns \$30B of mature profit while Segment~B posts a \$4B operating loss building early-stage R\&D options, applying a $20\times$ P/E multiple to consolidated earnings implicitly assigns $20 \times (-4) = -80$ (that is, -\$80B of negative value) to Segment~B, even though management can cap losses or spin Segment~B out at positive equity value. Always value mature cash engines and early option bets separately.
  \end{itemize}

  \item \textbf{Cost of Revenues and partner distribution share versus Operating Expenses.}
  GAAP separates \textbf{Cost of Revenues} (direct variable and capacity costs required to deliver each unit of revenue -- hosting, manufacturing, payment-network interchange, content royalties, and third-party distribution revenue share) from \textbf{Operating Expenses} (\textbf{R\&D}, \textbf{Sales \& Marketing}, and \textbf{General \& Administrative} period overhead). Whenever a platform relies on third-party channels (browsers, mobile OEMs, app stores, resellers) to route traffic or transactions, track partner revenue share as a percentage of channel revenue. A rising partner take-rate signals weakening distribution bargaining power and compresses Gross Margin before top-line revenue ever slows.

  \item \textbf{Stock-Based Compensation (SBC) and per-share cash dilution (ASC~718).}
  Employee equity grants (restricted stock units and stock options) are recorded as a GAAP compensation expense inside Cost of Revenues, R\&D, S\&M, and G\&A because equity is real pay: eliminating stock grants would force equivalent cash salary increases to retain talent. On the Cash Flow Statement, however, SBC is added back to Net Income to arrive at Operating Cash Flow ($\text{OCF}$) because issuing shares consumes zero cash on vesting day. Whenever management touts ``Non-GAAP EPS'' or raw Free Cash Flow ($\FCF = \text{OCF} - \text{CapEx}$) that ignores SBC, run two checks:
  \begin{align*}
    \FCF_{\text{SBC-burdened}} &= \FCF - \text{SBC}, \\
    \text{Net Share Retirement Cash} &= (\text{Share Buybacks} + \text{RSU Tax Withholdings}) - \text{SBC}.
  \end{align*}
  Only buyback cash spent \emph{above} quarterly SBC actually shrinks the diluted share count; if a company makes zero buybacks, inspect year-over-year growth in Diluted Shares Outstanding to measure the annual dilution tax on existing owners.

  \item \textbf{Capital Expenditures (CapEx) versus lagging GAAP Depreciation (D\&A).}
  Buying long-lived physical assets (factories, servers, GPU clusters, fiber networks) triggers an immediate $100\%$ cash outflow on Day~1 inside Investing Cash Flows (cutting $\FCF$ dollar for dollar), while recording \textbf{zero Day-1 expense} on the GAAP Income Statement. Instead, the asset sits on the Balance Sheet under Property, Plant \& Equipment (PP\&E) and is expensed gradually via straight-line \textbf{Depreciation} over its estimated useful life $L$ ($4$--$6$ years for servers; $15$--$25$ years for buildings). Consequently, whenever a company enters an investment surge where $\text{CapEx} / \text{D\&A} \gg 1$, current GAAP Net Income sits substantially above $\FCF$, and future GAAP operating margins face a predictable multi-year wave of rising depreciation charges. Notice also that changing the accounting estimate $L$ (say from $4$ years to $6$ years) raises GAAP Net Income immediately while leaving $\FCF$ completely unchanged.
\end{enumerate}

\subsection{Case study: hyperscaler quarterly release (Alphabet archetype)}

Table~\ref{tab:ch04-alphabet-archetype} applies all four principles to a representative single-quarter consolidated release modeled closely on Alphabet Inc.'s reporting architecture and scale (figures in USD billions, except per-share data).

\begin{table}[htbp]
\centering
\caption{Illustrative Case Study -- Hyperscaler Quarterly Release Excerpt (Alphabet Reporting Archetype, USD billions except EPS and share count).}
\label{tab:ch04-alphabet-archetype}
\small
\begin{tabular}{lrr}
\toprule
\textbf{Line Item} & \textbf{Prior Q3} & \textbf{Current Q3} \\
\midrule
\multicolumn{3}{l}{\emph{Panel A: Disaggregated Revenue}} \\
Google Search \& other                          & \$44.0 & \$49.4 \\
YouTube ads                                     & 8.0    & 8.9    \\
Google Network                                  & 7.7    & 7.5    \\
\quad \emph{Subtotal: Google Advertising}       & \emph{59.7} & \emph{65.8} \\
Google subscriptions, platforms, and devices    & 8.3    & 10.7   \\
\quad \emph{Subtotal: Google Services}          & \emph{68.0} & \emph{76.5} \\
Google Cloud                                    & 8.4    & 11.4   \\
Other Bets                                      & 0.3    & 0.4    \\
\textbf{Total Consolidated Revenues}            & \textbf{\$76.7} & \textbf{\$88.3} \\
\midrule
\multicolumn{3}{l}{\emph{Panel B: Consolidated GAAP Income Statement}} \\
Cost of revenues -- Traffic acquisition costs (TAC) & \$12.6 & \$13.7 \\
Cost of revenues -- Other (data centers, content, hardware) & 20.6 & 23.3 \\
\quad \emph{Total cost of revenues}             & \emph{33.2} & \emph{37.0} \\
\textbf{Gross Profit}                           & \textbf{\$43.5} & \textbf{\$51.3} \\
Research and development (R\&D)                 & 11.3   & 12.4   \\
Sales and marketing (S\&M)                      & 6.9    & 7.2    \\
General and administrative (G\&A)               & 4.0    & 3.2    \\
\quad \emph{Total operating expenses}           & \emph{22.2} & \emph{22.8} \\
\textbf{GAAP Operating Income}                  & \textbf{\$21.3} & \textbf{\$28.5} \\
Other income (expense), net                     & 0.8    & 3.2    \\
Income before income taxes                      & 22.1   & 31.7   \\
Provision for income taxes                      & 2.4    & 5.4    \\
\textbf{GAAP Net Income}                        & \textbf{\$19.7} & \textbf{\$26.3} \\
Diluted shares outstanding (billions)           & 12.68  & 12.45  \\
\textbf{Diluted GAAP EPS}                       & \textbf{\$1.55} & \textbf{\$2.11} \\
\midrule
\multicolumn{3}{l}{\emph{Panel C: Segment Operating Income (Loss)}} \\
Google Services                                 & \$23.9 & \$30.9 \\
Google Cloud                                    & 0.3    & 1.9    \\
Other Bets                                      & (1.2)  & (1.1)  \\
Unallocated corporate costs                     & (1.7)  & (3.2)  \\
\textbf{Total GAAP Operating Income}            & \textbf{\$21.3} & \textbf{\$28.5} \\
\midrule
\multicolumn{3}{l}{\emph{Panel D: Cash Flow \& Share Repurchase Excerpt (Current Q3)}} \\
GAAP Net Income                                 &        & \$26.3 \\
Depreciation and intangible amortization (D\&A) &        & 4.0    \\
Stock-based compensation (SBC) expense add-back &        & 5.8    \\
Working capital, deferred tax, and other operating items & & (5.4)  \\
\textbf{Net Cash Provided by Operating Activities (OCF)} & & \textbf{\$30.7} \\
Less: Purchases of property and equipment (CapEx) &      & (13.1) \\
\textbf{Free Cash Flow ($\FCF$)}                &        & \textbf{\$17.6} \\
\emph{Memo: Share repurchases (\$15.3B) + RSU tax settlements (\$2.4B)} & & \emph{(17.7)} \\
\bottomrule
\end{tabular}
\end{table}

Walk through our four diagnostic checks on Table~\ref{tab:ch04-alphabet-archetype} and Figure~\ref{fig:ch04-hyperscaler-bridge}:
\begin{itemize}
  \item \textbf{Segment physics:} Consolidated revenue rose $15.1\%$ YoY (\$76.7B to \$88.3B), whereas GAAP Operating Income rose $33.8\%$ (\$21.3B to \$28.5B). Panel~C reveals why: Google Services generated \$30.9B ($40.4\%$ segment margin), Google Cloud crossed fixed infrastructure scale -- adding $\Delta R =$ \$3.0B of revenue (\$8.4B to \$11.4B) and $\Delta \text{OI} =$ \$1.6B of profit (\$0.3B to \$1.9B) for a $53.3\%$ incremental operating margin -- while Other Bets (autonomous driving and life-science ventures) posted a \$1.1B operating loss, requiring standalone sum-of-the-parts treatment.
  \item \textbf{Distribution moat check (TAC):} \textbf{Traffic Acquisition Costs} (\$13.7B) measure revenue-share payments to OEMs and browsers for default search placement plus third-party Network publishers. Dividing TAC by Google Advertising revenue gives a \textbf{TAC Rate} of $13.7 / 65.8 = 20.8\%$, down from $12.6 / 59.7 = 21.1\%$ a year earlier as first-party Search and YouTube outgrew higher-TAC Network ads.
  \item \textbf{SBC and buyback check:} Headline $\FCF$ is \$17.6B, while $\FCF_{\text{SBC-burdened}} = 17.6 - 5.8 =$ \$11.8B. Because management deployed \$15.3B in buybacks plus \$2.4B in employee RSU tax settlements (\$17.7B total equity cash outflow), \$5.8B offset quarterly SBC grants and the remaining \$11.9B retired shares net, shrinking diluted share count by $1.8\%$ YoY ($12.68$B to $12.45$B).
  \item \textbf{AI CapEx versus lagging D\&A:} Quarterly CapEx for TPUs/GPUs and data centers hit \$13.1B against just \$4.0B of D\&A ($\text{CapEx} / \text{D\&A} = 3.28\times$). That \$9.1B excess explains why $\FCF$ (\$17.6B) sits \$8.7B below GAAP Net Income (\$26.3B) today, and signals that quarterly GAAP depreciation inside Cost of Revenues and R\&D will step up by billions of dollars over the coming four to six years.
\end{itemize}

\begin{figure}[htbp]
\centering
\begin{tikzpicture}[x=1.52cm, y=0.052cm]
  % Baseline and y-axis
  \draw[gray!50] (0.35,0) -- (7.75,0);
  \draw[->, >=stealth, thick, gray!80!black] (0.35,0) -- (0.35,96) node[above, font=\scriptsize\bfseries, text=black] {USD billions};
  \foreach \y in {0,20,40,60,80} {
    \draw[gray!25, dashed] (0.35,\y) -- (7.75,\y);
    \node[left, font=\scriptsize, text=gray!80!black] at (0.33,\y) {\y};
  }

  % Bar 1: Revenue (0 to 88.3)
  \filldraw[fill=sourceblue!75, draw=sourceblue!95!black] (0.65,0) rectangle (1.25,88.3);
  \node[above, font=\scriptsize\bfseries, text=sourceblue!95!black] at (0.95,88.3) {88.3};

  % Connector 1->2
  \draw[gray!70, densely dotted] (1.25,88.3) -- (1.65,88.3);
  % Bar 2: -TAC (88.3 down to 74.6, i.e. -13.7)
  \filldraw[fill=trapred!65, draw=trapred!90!black] (1.65,74.6) rectangle (2.25,88.3);
  \node[above, font=\scriptsize, text=trapred!90!black] at (1.95,88.3) {$-13.7$};

  % Connector 2->3
  \draw[gray!70, densely dotted] (2.25,74.6) -- (2.65,74.6);
  % Bar 3: -Other COGS (74.6 down to 51.3, i.e. -23.3)
  \filldraw[fill=trapred!65, draw=trapred!90!black] (2.65,51.3) rectangle (3.25,74.6);
  \node[below, font=\scriptsize, text=trapred!90!black] at (2.95,51.0) {$-23.3$};

  % Connector 3->4
  \draw[gray!70, densely dotted] (3.25,51.3) -- (3.65,51.3);
  % Bar 4: -OpEx (51.3 down to 28.5, i.e. -22.8)
  \filldraw[fill=trapred!65, draw=trapred!90!black] (3.65,28.5) rectangle (4.25,51.3);
  \node[above, font=\scriptsize, text=trapred!90!black] at (3.95,51.3) {$-22.8$};

  % Connector 4->5 (Operating Income 28.5 + Other 3.2 - Tax 5.4 = Net Income 26.3)
  \draw[gray!70, densely dotted] (4.25,28.5) -- (4.65,28.5);
  % Bar 5: GAAP Net Income total bar (0 to 26.3, showing -2.2 net non-op/tax step)
  \filldraw[fill=structurepurple!65, draw=structurepurple!90!black] (4.65,0) rectangle (5.25,26.3);
  \filldraw[fill=trapred!45, draw=trapred!80!black] (4.65,26.3) rectangle (5.25,28.5);
  \node[above, font=\scriptsize\bfseries, text=structurepurple!95!black] at (4.95,28.5) {26.3};

  % Connector 5->6 (Net Inc 26.3 + D&A 4.0 + SBC 5.8 - WC 5.4 = OCF 30.7, net +4.4)
  \draw[gray!70, densely dotted] (5.25,26.3) -- (5.65,26.3);
  % Bar 6: Operating Cash Flow (0 to 30.7)
  \filldraw[fill=quotientorange!55, draw=quotientorange!90!black] (5.65,0) rectangle (6.25,26.3);
  \filldraw[fill=targetgreen!55, draw=targetgreen!85!black] (5.65,26.3) rectangle (6.25,30.7);
  \node[above, font=\scriptsize\bfseries, text=quotientorange!90!black] at (5.95,30.7) {30.7};

  % Connector 6->7 (OCF 30.7 - CapEx 13.1 = FCF 17.6)
  \draw[gray!70, densely dotted] (6.25,30.7) -- (6.65,30.7);
  % Bar 7: CapEx drop (30.7 down to 17.6) and FCF bar (0 to 17.6)
  \filldraw[fill=trapred!55, draw=trapred!85!black] (6.65,17.6) rectangle (7.25,30.7);
  \filldraw[fill=targetgreen!75, draw=targetgreen!95!black] (6.65,0) rectangle (7.25,17.6);
  \node[above, font=\scriptsize, text=trapred!90!black] at (6.95,30.7) {$-13.1$};
  \node[font=\scriptsize\bfseries, text=white] at (6.95,9.5) {17.6};

  % Dashed SBC-burdened FCF line on Bar 7 (11.8)
  \draw[white, very thick, dashed] (6.65,11.8) -- (7.25,11.8);
  \draw[<-, >=stealth, thin, black] (7.28,11.8) -- (7.62,14.5) node[above left, font=\tiny, fill=white, inner sep=1pt] {ex-SBC: 11.8};

  % X-axis labels
  \node[below, align=center, font=\tiny\bfseries] at (0.95,-1) {Rev.};
  \node[below, align=center, font=\tiny] at (1.95,-1) {TAC};
  \node[below, align=center, font=\tiny] at (2.95,-1) {Other\\COGS};
  \node[below, align=center, font=\tiny] at (3.95,-1) {OpEx\\(incl.\ SBC)};
  \node[below, align=center, font=\tiny\bfseries] at (4.95,-1) {GAAP\\Net Inc.};
  \node[below, align=center, font=\tiny\bfseries] at (5.95,-1) {Oper.\\Cash Flow};
  \node[below, align=center, font=\tiny\bfseries] at (6.95,-1) {CapEx \&\\$\FCF$};
\end{tikzpicture}
\caption{Quarterly Revenue-to-Free-Cash-Flow waterfall bridge (Table~\ref{tab:ch04-alphabet-archetype}, USD billions). Starting from \$88.3B of revenue, deducting \$13.7B of TAC, \$23.3B of other COGS, \$22.8B of OpEx, and \$2.2B of net tax and non-operating outflows yields \$26.3B of GAAP Net Income. Adding back \$4.0B of D\&A and \$5.8B of SBC (net of \$5.4B in working-capital and tax timing outflows) produces \$30.7B of Operating Cash Flow, from which \$13.1B of AI/cloud CapEx leaves \$17.6B of $\FCF$ (\$11.8B after deducting \$5.8B of employee stock grants).}
\label{fig:ch04-hyperscaler-bridge}
\end{figure}

\section{Unit Economics, Cohort Retention, and S-1 IPO Prospectuses}

High-growth technology companies filing a Form~S-1 (or raising late-stage venture capital) frequently report rapid revenue growth alongside large GAAP operating losses. How do you distinguish a durable compounding machine from a structurally broken subsidy business that loses money on every unit sold?

\subsection{Core principles: unit contribution versus fixed step-costs and cohort metrics}

Start by splitting all costs into \textbf{unit-level variable delivery costs} (captured in Gross Margin $\text{GM} = \frac{R - \text{COGS}}{R}$) and \textbf{period R\&D plus overhead costs} ($F_{\text{total}} = F_{\text{R\&D}} + \text{OpEx}_{\text{fixed}}$). Operating profit as a function of revenue $R$ is
\begin{equation}
  \Pi(R) = \text{GM} \cdot R - F_{\text{total}}.
  \label{eq:unit-vs-fixed}
\end{equation}
Two structural regimes emerge immediately from \eqref{eq:unit-vs-fixed}:
\begin{itemize}
  \item \textbf{Negative unit margin ($\text{GM} < 0$):} Each dollar of new revenue destroys $|\text{GM}|$ dollars before paying a single engineer or salesperson ($\dv{\Pi}{R} = \text{GM} < 0$). Scaling top-line volume accelerates bankruptcy unless unit delivery cost falls below price (the classic failure mode of subsidized movie-ticket subscriptions, below-cost instant grocery delivery, and underpriced ``unlimited'' API wrappers).
  \item \textbf{Positive unit margin ($\text{GM} > 0$) funding upfront R\&D step-costs:} Each incremental dollar of revenue contributes $+\text{GM}$ dollars toward covering $F_{\text{total}}$, yielding an exact operating break-even revenue threshold
  \begin{equation}
    R^* = \frac{F_{\text{total}}}{\text{GM}}.
    \label{eq:breakeven-scale}
  \end{equation}
  Always check the \emph{economic half-life} of what $F_{\text{R\&D}}$ builds. If R\&D builds sticky enterprise workflow software whose codebase generates revenue for ten years, $F_{\text{R\&D}}$ flattens out as a percentage of revenue. If R\&D builds an asset that becomes obsolete every $12$ to $18$ months, then $F_{\text{R\&D}}$ is a recurring generational step-cost, and long-run profitability requires gross profit dollars ($\text{GM} \cdot R$) to compound faster across generations than $F_{\text{R\&D}}$.
\end{itemize}

Next, because aggregate revenue growth can mask severe customer defection behind heavy marketing spend, measure six universal cohort and capital-efficiency metrics in every S-1:

\begin{enumerate}
  \item \textbf{Annualized Run-Rate Revenue ($\text{ARR}$):} For contracted annual subscriptions, $\text{ARR}$ equals the annualized value of active recurring contracts ($12 \times \text{contracted Monthly Recurring Revenue}$). When revenue comes from uncontracted pay-as-you-go usage, multiplying a peak month by $12$ overstates durable recurring demand; always separate \emph{committed contracted ARR} from \emph{annualized usage run-rate}.

  \item \textbf{Gross Dollar Retention ($\text{GDR}$) and Net Dollar Retention ($\text{NDR}$):} Fix the cohort of customers active at $t = 0$ with starting ARR $R_0$. One year later ($t = 1$), excluding all new customers signed during the year, measure lost ARR from full cancellations ($C_{\text{churn}}$), lost ARR from downgrades ($D_{\text{contr}}$), and expansion ARR from higher seats or usage inside surviving accounts ($E_{\text{exp}}$):
  \[
    \text{GDR} = \frac{R_0 - C_{\text{churn}} - D_{\text{contr}}}{R_0} \le 100\%, \qquad
    \text{NDR} = \frac{R_0 - C_{\text{churn}} - D_{\text{contr}} + E_{\text{exp}}}{R_0}.
  \]
  $\text{GDR}$ measures floor stickiness ($90\%$--$97\%$ in mission-critical enterprise infrastructure; often $45\%$--$65\%$ in experimental consumer tools). $\text{NDR} > 100\%$ means a cohort's dollar revenue compounds organically every year even with zero new customer additions.

  \item \textbf{Customer Acquisition Cost ($\text{CAC}$) and Gross-Margin-Adjusted Payback:} Dividing new-business Sales \& Marketing expense (lagged by the average sales cycle) by new customers acquired gives $\text{CAC}$. If a customer generates monthly revenue $\text{ARPU}_{\text{mo}}$ at gross margin $\text{GM}$, the months required to earn back $\text{CAC}$ in gross profit are
  \[
    \text{CAC Payback (months)} = \frac{\text{CAC}}{\text{ARPU}_{\text{mo}} \times \text{GM}}.
  \]
  Never omit $\text{GM}$ from the denominator: a $55\%$-margin product takes $1.45\times$ longer to recoup $\text{CAC}$ than an $80\%$-margin product at the same price.

  \item \textbf{Lifetime Value to CAC ($\text{LTV} / \text{CAC}$):} With constant monthly net churn rate $c_{\text{mo}} \in (0, 1]$, expected lifetime is $\sum_{k=0}^\infty (1 - c_{\text{mo}})^k = 1 / c_{\text{mo}}$ months, giving undiscounted gross-profit $\text{LTV}$ and acquisition multiple
  \[
    \text{LTV} = \frac{\text{ARPU}_{\text{mo}} \times \text{GM}}{c_{\text{mo}}}, \qquad
    \frac{\text{LTV}}{\text{CAC}} = \frac{\text{ARPU}_{\text{mo}} \times \text{GM}}{c_{\text{mo}} \cdot \text{CAC}} = \frac{1}{c_{\text{mo}} \cdot \text{CAC Payback (months)}}.
  \]
  Healthy enterprise economics typically clear $\text{LTV}/\text{CAC} \ge 3.0\times$ and payback $\le 18$ months.

  \item \textbf{Burn Multiple:} To capture cash efficiency across the \emph{entire} company (including R\&D and G\&A, not just S\&M), divide net cash consumed by net new recurring revenue added:
  \[
    \text{Burn Multiple} = \frac{\text{Net Cash Burn in Period}}{\text{Net New ARR Added in Period}}.
  \]
  A Burn Multiple below $1.5\times$ indicates disciplined growth; above $2.5\times$ indicates that the company burns more than \$2.50 of cash for every \$1.00 of recurring top-line added.

  \item \textbf{The Rule of 40:} For scaled software and cloud platforms, evaluate growth-versus-profitability balance via
  \[
    \text{Rule of 40 Score} = \text{YoY Revenue Growth (\%)} + \text{Free Cash Flow Margin (\%)} \ge 40\%.
  \]
\end{enumerate}

\begin{trapbox}
\textbf{Four recurring red flags in S-1 filings and pre-IPO data rooms:}
\begin{enumerate}
  \item \textbf{``Community-Adjusted EBITDA'' and recurring-cost exclusions.} WeWork's 2019 S-1 invented ``Community-Adjusted EBITDA'' by stripping out building-level marketing, tenancy setup, design, and corporate G\&A alongside interest, taxes, D\&A, and SBC -- turning a \$1.9B GAAP loss into a fictional \$467M adjusted profit. Reject any metric that excludes ongoing costs of operating or acquiring customers.
  \item \textbf{Capitalizing routine software maintenance to inflate EBITDA.} Under ASC~350-40, qualifying internal-use software development can be capitalized onto the Balance Sheet rather than expensed in R\&D. Aggressive issuers reclassify routine maintenance as capitalized software, lifting Adjusted EBITDA while cash still exits on the Cash Flow Statement under Investing Activities (\emph{Capitalized internal-use software}). Always deduct capitalized software when computing $\FCF$.
  \item \textbf{Circular cloud credits and round-tripped barter ARR.} When a strategic cloud investor puts \$2.0B into a startup (or issues \$2.0B of promotional GPU compute credits) and the startup commits \$2.0B back to that investor's cloud, check the \emph{Related-Party Transactions} footnote: non-cash compute credits depress reported cash burn temporarily until credits expire and cash GPU bills hit COGS and R\&D in full. Likewise, watch for reciprocal startup deals where Startup~A buys \$5M of Startup~B's software in exchange for Startup~B buying \$5M of Startup~A's software.
  \item \textbf{Customer concentration.} Issuers must disclose any single customer representing $\ge 10\%$ of revenue. When $35\%$ of ARR rests on two fast-moving customers capable of insourcing the product, headline ARR carries hidden binary renewal risk.
\end{enumerate}
\end{trapbox}

\subsection{Case study: frontier AI lab economics (Anthropic archetype)}

Apply Equations~\eqref{eq:unit-vs-fixed}--\eqref{eq:breakeven-scale} and our cohort metrics to a representative frontier AI laboratory generating $R =$ \$3.20 billion of annualized API and enterprise seat revenue while reporting a \$1.44 billion GAAP operating loss:
\begin{itemize}
  \item \textbf{Unit serving (inference) margin:} Suppose traffic averages \$6.00 of revenue per million blended tokens ($80\%$ input at \$3.00/M and $20\%$ output at \$18.00/M). Direct serving costs per million tokens are \$0.60 in cloud marketplace revenue-share ($10.0\%$), \$1.35 in dedicated inference accelerator amortization ($22.5\%$), and \$0.45 in datacenter power, cooling, and networking ($7.5\%$) -- totaling \$2.40/M tokens ($40.0\%$). Thus \textbf{Inference Gross Profit} is \$3.60/M tokens, or $\text{GM}_{\text{inf}} = 60.0\%$. On $R =$ \$3.20B of revenue, inference COGS is \$1.28B and Inference Gross Profit is \$1.92B.
  \item \textbf{Generational pre-training step-cost:} Below gross profit, the lab spends $F_{\text{gen}} =$ \$2.70B training and aligning its next-generation frontier models (whose commercial pricing power lasts roughly $12$--$18$ months before rival releases commoditize that capability tier) plus $\text{OpEx}_{\text{fixed}} =$ \$0.66B on S\&M and G\&A ($F_{\text{total}} =$ \$3.36B), producing a GAAP operating loss of $1.92 - 3.36 =$ -\$1.44B.
  \item \textbf{Break-even scale and cohort durability:} By \eqref{eq:breakeven-scale}, current-generation operating break-even sits at $R^* = 3.36 / 0.60 =$ \$5.60B (Figure~\ref{fig:ch04-ai-cohorts}(a)). If enterprise cohorts exhibit $\text{GDR} = 92\%$ and $\text{NDR} = 125\%$ (losing \$8M of a \$100M starting cohort to churn/contraction while expanding \$33M in surviving accounts, Figure~\ref{fig:ch04-ai-cohorts}(b)), with $\text{ARPU}_{\text{mo}} =$ \$4,000, $\text{CAC} =$ \$36,000, and $c_{\text{mo}} = 1.0\%$, then gross-margin-adjusted CAC payback is $36{,}000 / (4{,}000 \times 0.60) = 15.0$ months and $\text{LTV} / \text{CAC} = 6.67\times$. Provided gross profit dollars ($0.60 R$) compound faster across generations than next-generation training clusters $F_{\text{gen}}$, the underlying serving business is a high-return compounder rather than a negative-margin subsidy furnace.
\end{itemize}

\begin{figure}[htbp]
\centering
\begin{tikzpicture}
  % Panel (a): Frontier AI Operating Profit vs Revenue Scale
  \begin{scope}[shift={(0,0)}]
    \begin{axis}[
      width=5.7cm, height=5.1cm,
      title={Panel (a): Serving Margin vs.\ Training Step},
      title style={font=\footnotesize\bfseries, text=sourceblue},
      xlabel={annual revenue $R$ (USD billions)},
      ylabel={operating profit (USD billions)},
      label style={font=\footnotesize},
      tick label style={font=\scriptsize},
      xmin=0, xmax=9.0, ymin=-4.2, ymax=2.4,
      axis lines=left,
      xtick={0,2,4,5.6,8}, ytick={-4,-2,0,2},
      enlargelimits=false, clip=false
    ]
      \draw[gray!60, thin] (axis cs:0,0) -- (axis cs:9.0,0);
      % Positive 60% inference GM minus USD 3.36B fixed pre-training + OpEx
      \addplot[targetgreen, very thick, domain=0:9.0] {0.60*x - 3.36};
      % Negative -40% inference GM wrapper minus USD 0.4B OpEx
      \addplot[trapred, very thick, dashed, domain=0:7.5] {-0.40*x - 0.40};
      % Current point R = 3.20, Profit = -1.44
      \filldraw[sourceblue] (axis cs:3.20,-1.44) circle (2pt);
      \node[sourceblue, font=\tiny, below right] at (axis cs:3.25,-1.48) {$R=3.2,\;\Pi=-1.44$};
      % Break-even R* = 5.60
      \filldraw[targetgreen!85!black] (axis cs:5.60,0) circle (2pt);
      \node[targetgreen!85!black, font=\tiny\bfseries, above left] at (axis cs:5.55,0.15) {$R^*=5.6$};
      \node[targetgreen!85!black, font=\tiny, right] at (axis cs:4.8,1.75) {$\text{GM}_{\text{inf}}=+60\%$};
      \node[trapred, font=\tiny, right] at (axis cs:2.8,-3.55) {$\text{GM}_{\text{inf}}=-40\%$ (subsidy)};
    \end{axis}
  \end{scope}
  % Panel (b): Cohort Dollar Retention trajectories
  \begin{scope}[shift={(6.4,0)}]
    \begin{axis}[
      width=5.7cm, height=5.1cm,
      title={Panel (b): \$100M Starting Cohort Over Time},
      title style={font=\footnotesize\bfseries, text=targetgreen},
      xlabel={years after cohort acquisition},
      ylabel={cohort ARR remaining (USD millions)},
      label style={font=\footnotesize},
      tick label style={font=\scriptsize},
      xmin=0, xmax=4.0, ymin=0, ymax=260,
      axis lines=left,
      xtick={0,1,2,3,4}, ytick={0,50,100,150,200,250},
      enlargelimits=false, clip=false
    ]
      % Enterprise AI/Cloud NDR = 125%
      \addplot[targetgreen, very thick, mark=*, mark size=1.5pt] coordinates {
        (0,100) (1,125) (2,156.25) (3,195.3) (4,244.1)
      };
      % Sticky SMB NDR = 100%
      \addplot[curveblue, very thick, dashed, mark=square*, mark size=1.3pt] coordinates {
        (0,100) (1,100) (2,100) (3,100) (4,100)
      };
      % High-churn wrapper NDR = 55%
      \addplot[trapred, very thick, densely dotted, mark=triangle*, mark size=1.6pt] coordinates {
        (0,100) (1,55) (2,30.25) (3,16.6) (4,9.15)
      };
      \node[targetgreen!85!black, font=\tiny\bfseries, left] at (axis cs:3.9,238) {$\text{NDR}=125\%$};
      \node[curveblue, font=\tiny\bfseries, above] at (axis cs:2.6,102) {$\text{NDR}=100\%$};
      \node[trapred, font=\tiny\bfseries, above right] at (axis cs:1.5,38) {$\text{NDR}=55\%$ (leaky)};
    \end{axis}
  \end{scope}
\end{tikzpicture}
\caption{Panel (a) contrasts a frontier AI lab with $+60\%$ inference gross margin and \$3.36B of generational training plus fixed OpEx (green, crossing operating break-even at $R^* =$ \$5.60B) against a negative-gross-margin subsidy business (red dashed, where scaling revenue deepens losses monotonically). Panel (b) tracks \$100M of starting cohort ARR over four years: at $\text{NDR} = 125\%$ the cohort compounds to \$244M with zero new sales, whereas at $\text{NDR} = 55\%$ over $90\%$ of the cohort evaporates.}
\label{fig:ch04-ai-cohorts}
\end{figure}

\section{Judging Profitability of a Small or Mid-Sized Private Business}

Evaluating an established private operating business doing \$2 million to \$50 million in revenue -- a specialty distributor, logistics fleet, B2B agency, HVAC contractor, or regional payments firm -- differs fundamentally from reading an SEC filing. Private books are rarely audited, often use cash or modified-cash accounting, and mingle founder compensation and tax-minimization choices with operating costs.

\subsection{Core principles: SDE versus Normalized EBITDA, proof of cash, and NWC pegs}

Brokers marketing founder-owned businesses commonly quote \textbf{Seller's Discretionary Earnings ($\text{SDE}$)}: pre-tax net income plus interest, depreciation, one working owner's total compensation, personal owner perks, and claimed ``one-time'' expenses. $\text{SDE}$ measures the combined return to \emph{both} your capital and $55$ hours per week of your own labor if you step in as full-time manager.

An investor or acquirer pricing the \emph{business} as a standalone asset must convert claimed $\text{SDE}$ into \textbf{True Normalized EBITDA} and \textbf{Normalized Pre-Tax Cash Flow} via five universal diligence rules:

\begin{enumerate}
  \item \textbf{Market-rate management replacement.} Add back current founder salaries ($+S_{\text{founder}}$) \emph{and simultaneously deduct} the fully loaded market salary required to hire qualified replacement executives ($-S_{\text{market}}$). When two founders take \$60,000 W-2 salaries (drawing remaining profit as distributions) while doing the work of a \$210,000 General Manager and a \$180,000 Sales Director, the net payroll adjustment to unadjusted EBITDA is $+120{,}000 - 390{,}000 =$ -\$270,000, rather than a \$120,000 positive add-back. (Conversely, if a passive owner draws \$600,000 while an existing salaried GM already runs operations, $+S_{\text{founder}}$ is a genuine positive adjustment.)
  \item \textbf{Invoice-level audit of personal and ``one-time'' add-backs.} Allow personal owner add-backs (vehicles, family travel, club dues) only when verified by specific ledger invoices and removable without hurting staff or clients. Reject ``one-time'' add-backs for bad debts, inventory write-downs, equipment repairs, or legal disputes whenever multi-year ledger history shows comparable charges recurring regularly.
  \item \textbf{Three-way proof of cash.} Because unaudited QuickBooks files can be edited retroactively without leaving an external audit trail, reconcile $24$ months of internal P\&L totals against (a) signed federal income tax returns (IRS Forms 1120, 1120-S, or 1065, where overstating profit incurs real cash tax under penalty of perjury) and (b) raw monthly bank statements pulled directly from the bank (verifying customer cash receipts net of inter-account transfers).
  \item \textbf{Maintenance CapEx deduction and the Net Working Capital ($\text{NWC}$) Peg.} Subtract annual physical Maintenance CapEx required to keep trucks, machinery, or facilities operating steady-state. In addition, specify a contractual \textbf{$\text{NWC}$ Peg} ($\text{NWC}_{\text{peg}}$, usually trailing 12-month average $\text{AR} + \text{Inventory} - \text{AP}$) in the purchase agreement so closing cash consideration adjusts via
  \[
    \text{Cash Paid at Close} = \text{Headline Purchase Price} + \bigl(\text{NWC}_{\text{close}} - \text{NWC}_{\text{peg}}\bigr).
  \]
  Without an $\text{NWC}$ peg in a ``cash-free, debt-free'' deal, a seller could liquidate receivables at a discount, halt inventory purchases, and delay supplier payables right before closing -- sweeping working capital out as ``cash'' and forcing the buyer to fund a massive Day-1 working-capital hole.
  \item \textbf{Counterparty and key-person transferability.} Audit change-of-control clauses in supplier/franchise agreements and verify whether top customer accounts depend personally on the departing founder.
\end{enumerate}

\subsection{Case study: recasting a \$9.0M private B2B distributor}

A business broker markets a \$9.0M-revenue industrial fasteners distributor for \$8.25M, advertising ``$5.0\times$ cash flow'' on \$1.65M of claimed $\text{SDE}$. Applying our five normalization rules to the company's Form~1120-S tax return, bank statements, and general ledger produces the diligence bridge below (in USD thousands):

\begin{center}
\small
\begin{tabular}{lr}
\toprule
\textbf{Private Diligence Bridge Line Item} & \textbf{USD Thousands} \\
\midrule
Tax-return Pre-Tax Net Income                                           & \$920 \\
$+$ Net Interest Expense (\$90k) and Tax Depreciation (\$160k)          & $+250$ \\
\textbf{Unadjusted Tax-Return EBITDA}                                   & \textbf{\$1,170} \\
$+$ Two co-founders' W-2 salaries (\$60k each, added back by broker)    & $+120$ \\
$+$ Verified owner personal vehicles, club dues, and personal travel    & $+110$ \\
$+$ Broker's claimed ``one-time'' legal (\$90k) + shrinkage/bad-debt (\$160k) & $+250$ \\
\textbf{Broker's Claimed Seller's Discretionary Earnings ($\text{SDE}$)} & \textbf{\$1,650} \\
\midrule
$-$ Disallow recurring annual inventory shrinkage and customer bad debt & $-160$ \\
$-$ Market-rate replacement General Manager (\$210k) + Sales Director (\$180k) & $-390$ \\
\textbf{True Normalized EBITDA}                                         & \textbf{\$1,100} \\
$-$ Average annual Maintenance CapEx (delivery vans, warehouse racking) & $-220$ \\
\textbf{Normalized Pre-Tax Unlevered Operating Cash Flow}               & \textbf{\$880} \\
\bottomrule
\end{tabular}
\end{center}

Disallowing recurring bad-debt/inventory charges and charging market-rate replacement management lowers sustainable EBITDA by one-third (from \$1.65M to \$1.10M), and deducting \$220k of maintenance CapEx leaves \$880k of pre-tax unlevered cash flow. At the \$8.25M asking price, the buyer was actually being asked to pay $7.5\times$ Normalized EBITDA and $9.4\times$ pre-tax post-CapEx cash flow -- nearly double the broker's $5.0\times$ pitch.

Part~III shifts from evaluating established operating businesses to financing new companies from zero -- deriving the exact math of cap tables, SAFE conversions, priced venture rounds, option-pool shuffles, and liquidation waterfalls.

\section*{Exercises}
\addcontentsline{toc}{section}{Exercises}

\begin{exercise}[Dissecting a hyperscaler quarterly release]
A public cloud-and-advertising company reports the following single-quarter figures (in USD billions): Search \& first-party ad revenue \$38.0B, Network partner ad revenue \$6.0B, Cloud revenue \$9.5B (up from \$7.0B in the prior-year quarter), Other Bets revenue \$0.5B; Traffic Acquisition Costs \$9.24B, Other cost of revenues \$18.26B; R\&D \$10.2B, S\&M \$5.8B, G\&A \$2.5B. Segment operating income is +\$24.0B for Services, +\$1.6B for Cloud (up from +\$0.2B in the prior-year quarter), -\$0.9B for Other Bets, and -\$2.7B unallocated corporate costs. On the Cash Flow Statement, GAAP Net Income is \$19.4B, D\&A is \$3.5B, Stock-Based Compensation is \$4.6B, working-capital adjustments are -\$1.5B, CapEx is \$11.2B, and open-market share repurchases are \$12.0B.
(a) Verify Total Consolidated Revenues, Gross Profit, Gross Margin, and GAAP Operating Income, and confirm that the segment operating results reconcile to consolidated GAAP Operating Income.
(b) Compute the Advertising TAC Rate and Cloud's year-over-year incremental operating margin ($\Delta \text{OI} / \Delta \text{Rev}$).
(c) Compute Operating Cash Flow, GAAP Free Cash Flow ($\FCF$), SBC-burdened $\FCF$, and the dollar amount of share repurchases that genuinely reduced share count beyond offsetting quarterly SBC dilution.
\end{exercise}

\begin{exercise}[AI CapEx capitalization versus useful-life accounting]
On January~1 of Year~1, a cloud provider pays \$12.0 billion in cash to deploy a new fleet of AI server racks with zero salvage value.
(a) Assuming straight-line depreciation over a $4$-year useful life (and ignoring tax effects for simplicity), write down the exact dollar impact of this single \$12.0B purchase on Operating Cash Flow, Investing Cash Flow, Free Cash Flow, GAAP Depreciation Expense, and GAAP Operating Income in Year~1 versus Years~2--4.
(b) Suppose on January~1 of Year~1 management revises its accounting policy to depreciate AI servers over $6$ years instead of $4$ years. By how many billions of dollars does Year-1 GAAP Operating Income increase, and by how many dollars does Year-1 Free Cash Flow change?
\end{exercise}

\begin{exercise}[Frontier AI lab serving economics and break-even scale]
Over a fiscal year, a frontier AI lab bills $500$ million ``million-token blocks'' ($5.0 \times 10^{14}$ total tokens) across enterprise and API customers at a blended price of \$5.00 per million tokens. Direct serving costs per million tokens are \$0.40 in cloud partner fees, \$1.25 in inference GPU amortization, and \$0.35 in power and networking. During the same year, the lab spends \$1.95 billion on pre-training compute and research staff for its next-generation model and \$0.45 billion on S\&M and G\&A.
(a) Compute annual Revenue, Inference Cost of Revenue, Inference Gross Profit, Inference Gross Margin (\%), and GAAP Operating Loss.
(b) Calculate the annual revenue break-even $R^*$ required to cover the current \$2.40B fixed pre-training plus OpEx envelope at the current Inference Gross Margin.
(c) Suppose rival open-weights releases force the lab to cut blended API prices from \$5.00 to \$3.20 per million tokens while unit serving costs remain \$2.00 per million tokens, and next year's pre-training plus OpEx rises to \$3.00B. Compute the new Inference Gross Margin (\%) and the new break-even revenue $R^*$.
\end{exercise}

\begin{exercise}[S-1 cohort metrics and circular credits]
An enterprise AI software S-1 discloses the following metrics:
\begin{itemize}
  \item Customers active one year ago generated $R_0 =$ \$80.0M of ARR. Over the past twelve months, \$4.8M of that ARR churned completely, \$2.4M contracted via seat reductions, and surviving accounts expanded usage by \$23.2M.
  \item Including new logos, total ending ARR rose from \$80.0M to \$130.0M while reported Net Cash Burn for the year was \$65.0M (after utilizing \$20.0M of non-cash promotional cloud compute credits granted by its lead Series~C investor, which expire tomorrow).
  \item An average new enterprise customer pays $\text{ARPU}_{\text{mo}} =$ \$5,000/month at a true cash Gross Margin of $64\%$, fully loaded $\text{CAC}$ is \$44,800 per customer, and monthly net churn is $c_{\text{mo}} = 0.8\%$.
  \item YoY revenue growth is $62\%$ and true cash $\FCF$ margin (excluding promotional credits) is $-65\%$.
\end{itemize}
Compute (a) Gross Dollar Retention ($\text{GDR}$) and Net Dollar Retention ($\text{NDR}$), (b) Gross-Margin-Adjusted CAC Payback (in months) and undiscounted $\text{LTV} / \text{CAC}$, and (c) the headline Burn Multiple versus the true cash Burn Multiple (adjusting for the expiring \$20.0M investor cloud credits), alongside the Rule of 40 score.
\end{exercise}

\begin{exercise}[Private-company diligence and Net Working Capital peg]
You sign a letter of intent to acquire a private specialty testing laboratory at a headline valuation of $6.0\times$ Normalized EBITDA, cash-free and debt-free, subject to a contractual Net Working Capital ($\text{NWC}$) peg of \$1.40 million. The seller's QuickBooks P\&L reports \$1.15 million of pre-tax income after deducting \$80,000 of interest, \$140,000 of depreciation, a \$90,000 salary to the sole founder-CEO (who works full-time managing lab operations and client renewals), \$75,000 of verified personal boat and travel expenses, and \$105,000 of chemical waste-disposal fines that also appeared at \$95,000--\$110,000 in each of the prior three years. Hiring a qualified replacement Lab Director costs \$240,000/year fully loaded, and annual maintenance CapEx for spectrometer calibration and replacement is \$190,000.
(a) Compute the seller's claimed $\text{SDE}$ (assuming the seller adds back interest, depreciation, founder salary, personal perks, and the waste-disposal fines) versus your True Normalized EBITDA and Normalized Pre-Tax Cash Flow after Maintenance CapEx.
(b) At $6.0\times$ True Normalized EBITDA, what is the baseline Enterprise Value? If on closing day the seller delivers \$950,000 of Accounts Receivable, \$300,000 of Inventory, and \$410,000 of Accounts Payable, compute actual closing $\text{NWC}$ and the exact adjusted cash purchase price wired to the seller (assuming \$0 debt and \$0 retained cash).
\end{exercise}
